\chapter{Prediction Task, Data, and Conditions}
\label{ch:task}
% 개요 3장. 3.1, 3.2절은 기존 원고를 살리고 보강했다. 기존 3.3절(대상
% 식별과 정상 작업)은 대상 식별 부분을 3.2절의 UNKNOWN 문단으로 옮기고,
% 3.3절은 새 조건 정의로 바꿨다. 근거: RECORD:32--56,
% ORIG:PROPHET_ENV_CONSENSUS_KD_V3_20261120/PAPER_DEFINITIONS.md §1.

\section{An execution and a prediction row}
\label{sec:task-rows}

An execution is one run of a source procedure on the instrumented
Windows host. Each time the procedure requests its next action (a
\emph{source request}), the collection records a decision boundary
$t$, which also serves as the prediction cutoff. The task is to predict,
at each boundary, which of 18 semantic attacker actions comes next,
using only the Windows execution and audit events that the collector
received before $t$. The corresponding prediction row consists of that
evidence and the label of the next attacker action defined in
Chapter~\ref{ch:introduction}. The label is checked against the
action's own process creation record, and its termination record
supports verification when labels are constructed after the run.
Because one execution yields several boundaries, the number of rows
differs from the number of independent runs.
% 근거: RECORD:32. 라벨 정의 문장은 기존 원고 그대로 유지.

The 18-action vocabulary covers host, directory-service and network
enumeration, registry queries, and operations on synthetic files.
Appendix~\ref{app:actions} lists the categories. A process start or
termination is an observation, whereas the label identifies the
operation associated with the source request; the collector may
receive many other observations between two labeled actions.

\section{Receipt-bounded evidence}
\label{sec:receipt}

Let $r(e)$ denote the collector's receipt time for an event $e$,
and $\prec_r$ its order in the received stream. The input at boundary
$t$ is
\begin{equation}
  X_t = \operatorname{sort}_{\prec_r}\bigl\{e : r(e) < t\bigr\}.
  \label{eq:receipt}
\end{equation}
All admitted events precede the boundary strictly. When the collector
receives several events in one delivery (a \emph{receipt batch}), the
batch is atomic: either all of its events are admitted or none are, so
a boundary never falls inside a batch. The event-level GRU of the
baseline predictor (Chapter~\ref{ch:method}) nevertheless advances once
for each admitted event, which separates the collector's unit of
availability from the model's unit of computation.
% 근거: RECORD:51 (receipt < cutoff, 원자 receipt batch).

Windows Security event 4688 records process creation and, with
command-line auditing enabled, can include the command
line~\cite{ms4688}. Event 4689 records termination and includes an
application-specific exit status~\cite{ms4689}. We write START and END
for received 4688 and 4689 records. A START may be received before
its END, and at such a boundary the result of that process remains
unobserved. For target-aware processing, target identities are
recovered from observed command paths, and the same canonical target
can appear in several processes; an observed copy operation is given
distinct source and destination roles. When a received END has a unique
compatible START, it inherits the observed target relation of that
START. An END without such a START, or any record whose target cannot
be determined from the received evidence, is kept in the stream in an
\emph{UNKNOWN} state: it is neither dropped nor completed from the audit
ledger. The grouping by target is rebuilt at each boundary, and a prefix
that contains no usable file target is still mapped to 18 probabilities
through a target-free path.
% 근거: RECORD:51 (UNKNOWN 유지). UNKNOWN의 예: C-fold FIT의 native
% 기록 중 END 322건은 고유 START를 결속하지 못해 UNKNOWN으로 남았다
% (ORIG:PROPHET_SEED_INPUT_AUDIT_20261118/MECHANISM_BOUNDS_KO.md:28).
% TODO(사람): UNKNOWN이 "결속되지 않은 END/대상 미배정 record"만 뜻하는지,
% 아니면 END가 아직 수신되지 않은 START의 결과 상태도 포함하는지 확인해 주세요.

The predictor reads observed fields: the executable, available command
content, receipt order, elapsed time, and received exit status. Four
kinds of information are excluded from inference input and used only
for labeling and audit: the label itself, the execution identifier,
whether a record belongs to the source procedure or to a benign worker,
and any record received at or after the boundary. Process and record
identifiers are used to link START and END records and to audit the
collection.
% 근거: RECORD:50.

\begin{figure}[tb]
\centering
\begin{tikzpicture}[x=1cm,y=1cm,>=Stealth,
  event/.style={circle,inner sep=1.8pt,fill=blue!60!black},
  later/.style={circle,inner sep=1.8pt,fill=black!55},
  txt/.style={font=\footnotesize,align=center}]
  \fill[blue!5,rounded corners=2pt] (0.5,0.75) rectangle (7.35,1.65);
  \draw[->,black!64] (0.5,2.8) -- (12.1,2.8);
  \draw[->,black!64] (0.5,1.2) -- (12.1,1.2);
  \node[txt,anchor=west] at (0.5,3.6) {Windows event time};
  \node[txt,anchor=west] at (0.5,0.45) {Collector receipt time};
  \draw[thick,orange!75!black,dashed] (7.45,0.85) -- (7.45,3.65);
  \node[txt,text=orange!70!black] at (7.45,3.87) {Prediction boundary};
  \draw[->,blue!44!black!45] (2.15,2.74) -- (3.05,1.25);
  \draw[->,blue!44!black!45] (4.05,2.74) -- (5.0,1.25);
  \draw[->,black!48] (6.12,2.74) -- (8.78,1.25);
  \node[event,label={[txt]above:{P$_1$ starts}}] at (2.15,2.8) {};
  \node[event,label={[txt]above:{P$_2$ starts}}] at (4.05,2.8) {};
  \node[later,label={[txt]above:{P$_1$ exits}}] at (6.12,2.8) {};
  \node[later,label={[txt]above:{Next labeled\\action starts}}] at (10.62,2.8) {};
  \node[event,label={[txt]below:{P$_1$ START}}] at (3.05,1.2) {};
  \node[event,label={[txt]below:{P$_2$ START}}] at (5.0,1.2) {};
  \node[later,label={[txt]below:{P$_1$ END}}] at (8.78,1.2) {};
  \node[txt,anchor=east,text=blue!60!black] at (7.14,0.45) {Admitted prefix};
  \node[txt,anchor=west,text=black!65] at (8.12,0.45) {Received after boundary};
\end{tikzpicture}
\caption{A schematic prediction boundary. P$_1$ and P$_2$ denote two
observed processes. The predictor receives both START records before
the boundary. The END of P$_1$ arrives afterward even though P$_1$ has
already exited. The next attacker action supplies the evaluation label.}
\label{fig:cutoff}
\end{figure}

Figure~\ref{fig:cutoff} shows why availability cannot be reconstructed
from occurrence time alone: P$_1$ has already exited, yet its result is
not part of the admitted prefix. A later receipt can change later
forecasts but leaves the earlier prefix intact.

\section{Benign workloads and evaluation conditions}
\label{sec:conditions}

Some executions run benign tasks on the same host while the source
procedure executes. Three benign workers are defined. Worker~A repeats a
read-only \texttt{where} search, a \texttt{tasklist} process listing, and
an \texttt{sc} service listing every 1,200~ms in a separate scratch
directory. Worker~B repeats a \texttt{dir /b} directory listing and a
\texttt{certutil} SHA-256 hash of a scratch file at a fixed 1,200~ms
interval. Worker~C lists services with \texttt{sc} and then the local
interfaces with \texttt{ipconfig /all} at a fixed 2,400~ms interval; it
makes no external connection. A worker starts before the source is
dispatched and stops after the source finishes. It does not read the
source's actions, results, or requests, nor the defender's receipts,
and the source keeps its own result-dependent choice of actions. The
benign processes therefore enter the received stream without taking
part in the source's decisions. We call a condition with at least one
benign worker a \emph{background condition}.
Table~\ref{tab:conditions} lists the conditions and the two placements
of the file targets.
% 근거: RECORD:34--42; PAPER_DEFINITIONS.md §1.

\begin{table}[htbp]
\centering
\caption{Evaluation conditions and target layouts. ``C-fold FIT'' states
whether the condition's benign work occurs in the 603 training runs of
the C-fold (Section~\ref{sec:conditions}).}
\label{tab:conditions}
\begin{tabularx}{\textwidth}{@{}llX@{}}
\toprule
Identifier & C-fold FIT & Definition \\
\midrule
\texttt{ORIGINAL\_18} & seen & Original source procedure without a benign worker; all 18 actions \\
\texttt{C\_ONLY} & unseen & Source procedure with worker C \\
\texttt{A+C} & unseen & Source procedure with workers A and C \\
\texttt{B+C} & unseen & Source procedure with workers B and C \\
\midrule
\texttt{FLAT} & --- & Four file targets \texttt{echo}, \texttt{foxtrot}, \texttt{golf}, \texttt{hotel} placed as sibling directories \\
\texttt{NESTED} & --- & As \texttt{FLAT}, except that \texttt{foxtrot} is moved under \texttt{echo}; \texttt{echo} remains a target, and the two payloads have different search roots \\
\bottomrule
\end{tabularx}
\end{table}
% 표 기호: 본문·표의 조건 이름은 A+C FLAT, A+C NESTED, C_ONLY, B+C,
% ORIGINAL_18로 통일한다(개요 전체 규칙). 원고의 "C-only"는 C_ONLY로 바꿨다.
% TODO(사람): 배치 FLAT/NESTED는 A+C에서만 따로 보고됩니다. C_ONLY와 B+C
% 실행은 어느 배치를 쓰는지, 그리고 FIT 603실행에 NESTED 배치가 들어 있는지
% 알려 주세요. 표의 "---"를 채울 수 있습니다.

All new experiments in this thesis train on the C-fold, the outer fold
in which every variant of worker~C is withheld from the training runs.
We call these 603 runs the \emph{FIT set}. A condition is \emph{seen}
if its benign work occurs in the FIT set and \emph{unseen} otherwise;
under the C-fold, every condition that contains worker~C is unseen, and
\texttt{ORIGINAL\_18} is the seen reference. The B-fold withholds
worker~B instead. It appears only in the older comparisons of
Appendix~\ref{app:secondary-results} and is not used for any result in
the main chapters.
% 근거: RECORD:44; answers.md 7-(2). 이 정의로 05_evaluation.tex의 기존 TODO
% ("define the B-fold and C-fold")의 fold 부분을 닫는다. 그 TODO의 natural
% 전달 일정 정의는 [직접 기입]으로 남는다. 05장을 쓸 때 TODO를 정리해 주세요.

The C-fold makes the difference between seen and unseen conditions
concrete. Worker~C runs \texttt{ipconfig /all}, and the input
represents such records as an IP hint. Of the 4,824 boundaries in the
FIT set, 2,425 had received at least one IP hint. These hints came from
658 distinct native records in 329 runs, and all 658 belong to the
process lineage of the attack procedure. In the 60 \texttt{A+C}
development runs (660 rows), all 353 native records carrying an IP hint
belong instead to the lineage of benign worker~C. The same kind of
evidence is thus produced by the attacker during training and by benign
work at evaluation. Lineage was determined from audit information that
the model never receives. We use this case to show what a condition
unseen in training looks like; the audit does not establish it as the
cause of any seed's success or failure.
% 근거: ORIG:PROPHET_REL_GLOBAL_SEED_POPULATION_20261118/FINAL_REPORT_KO.md:9.
% 60실행·660행 = A+C FLAT 31·341 + NESTED 29·319 (같은 파일, 135실행 개발 집합).
% 4,824는 원 파일의 "자연 경계" 수다. natural boundary의 정의는 5장 [직접 기입].
% TODO(사람): "IP hint"가 ipconfig 사건에서 만든 입력 표현이라는 설명이 맞는지
% 확인해 주세요(원 파일은 FIT의 ipconfig.exe START/END 880개가 모두 공격
% 계보라고 적는다).

\section{A recorded file-processing decision}
\label{sec:recorded-case}

Table~\ref{tab:recorded-case} shows part of the received prefix for
execution \texttt{THME0064}, from the three-target development set,
before its first file-processing decision (\texttt{SEQ\_5}). The path
names distinguish three synthetic targets. The marker concerns
\texttt{alpha}, whose earlier payload search returned a nonzero exit
status, while the marker's own zero exit status belongs to a different
process. Within the \texttt{alpha} group, the latest received result is
therefore a success, and the search result comes from an earlier record.
Grouping the records by target lets a model read the two processes
together while keeping their results apart. The next recorded action
was a synthetic file search. The example previews the mechanism examined
in Chapter~\ref{ch:diagnosis}, where the in-house seed instability is
traced to the lookup of the current target's earlier result.

\begin{table}[htbp]
\centering
\caption{Selected records received before the decision in
\texttt{THME0064}. Rows show the observed target groups. For the
two other targets, only the search termination is shown.}
\label{tab:recorded-case}
\begin{tabularx}{\textwidth}{@{}llX@{}}
\toprule
Target group & Record ID & Evidence available before the decision \\
\midrule
\texttt{bravo}, \texttt{charlie} & & Payload search END at each target, zero exit status \\
\texttt{alpha} & 207481 & Payload search END, nonzero exit status \\
\texttt{alpha} & 207483 & Marker START \\
\texttt{alpha} & 207485 & Marker END, zero exit status \\
\bottomrule
\end{tabularx}
\end{table}
% 기존 원고의 target-wise GRU 대 relation attention 확률 비교 문장
% (0.434/0.450/0.481, 0.984/0.483/0.517)은 개요에 따라 부록 B.5로 옮긴다.
% 부록 B를 쓸 때 그 문장과 함께 "target-wise GRU: [직접 기입]" 자리를 남긴다.
% 기존 sec:membership 참조는 이 장에서 뺐다(소속 교환 설명은 부록 B에서 한 번만).

% ---------------------------------------------------------------
% 사람에게 물을 것 (3장)
% 1. UNKNOWN 상태의 범위: 결속되지 않은 END·대상 미배정 record만인지,
%    END가 아직 수신되지 않은 START의 결과도 UNKNOWN으로 부르는지.
% 2. FLAT/NESTED 배치가 C_ONLY, B+C 실행과 FIT 603실행에 어떻게 들어가는지
%    (표 tab:conditions의 "---" 칸).
% 3. "IP hint"의 정의: ipconfig 사건에서 만든 입력 표현이 맞는지.
% 4. THME0064가 속한 3대상 개발 집합이 C-fold FIT로 학습한 모델로 평가된
%    것인지(6장 기제 예고 문장의 범위 확인).
% ---------------------------------------------------------------
